%% kapa: exam=1 round=31 year=2020 period=2 form=A subject=appraisal_law subject_ko=감정평가관계법규 seq=39 number=39
%% source: Q-Net 감정평가사 기출문제 (cst00309) artlSeq=5207219 file=extracted/1008/5207219/A형/A형/2020년도 제31회 감정평가사 1차 시험_2교시_A형.pdf
%% conversion: kordoc 4.24.51 -> tools/gen1.py (auto); review=proofread-sonnet; flags=none
\begin{question}{39}
도시 및 주거환경정비법령상 조합설립추진위원회와 조합에 관한 설명으로 옳지 않은 것은?
\begin{choices}
  \item 조합설립추진위원회는 설계자의 선정 및 변경의 업무를 수행할 수 있다.
  \item 조합설립추진위원회는 추진위원회를 대표하는 추진위원장 1명과 감사를 두어야 한다.
  \item 조합장이 자기를 위하여 조합과 계약이나 소송을 할 때에는 이사가 조합을 대표한다.
  \item 정비사업전문관리업자의 선정 및 변경의 사항은 조합 총회의 의결을 거쳐야 한다.
  \item 조합장이 아닌 조합임원은 조합의 대의원이 될 수 없다.
\end{choices}
\end{question}
\answer{3}
